\documentclass[12pt]{article}
\usepackage[utf8]{inputenc}
\usepackage[spanish]{babel}
\usepackage{amsmath}
\usepackage{amsfonts}
\usepackage{amssymb}
\usepackage{geometry}
\geometry{letterpaper, margin=2.5cm}

\begin{document}

\begin{center}
    \Large \textbf{Evaluación de Examen 1: Unidad I} \\
    \large Física para Ciencias de la Salud (005-1713) \\
    \vspace{0.5cm}
    \textbf{Estudiante:} Karla Rodriguez \quad \textbf{C.I.:} 33.628.014
    \vspace{0.3cm} \\
    \large \textbf{Calificación Final: 9/10 puntos}
\end{center}

\vspace{0.5cm}

\section*{Análisis de Resultados}

\subsection*{1. Ángulo plano (Conversión a radianes)}
\textbf{Procedimiento y resultado:} Plantea correctamente la conversión usando el factor equivalente a $180^\circ$. Realiza una simplificación de fracciones de manera progresiva y muy limpia $\left(\frac{75\pi}{180} = \frac{15}{36} = \frac{5\pi}{12}\right)$, obteniendo el valor exacto de $\frac{5\pi}{12}$ rad. \\
\textbf{Veredicto:} \textbf{Correcto} (2/2 pts).

\subsection*{2. Sistemas de coordenadas (Longitud del hueso)}
\textbf{Procedimiento y resultado:} Utiliza la ecuación formal de la distancia entre dos puntos. Calcula por separado las diferencias en los ejes y las eleva al cuadrado. Aunque su notación es algo informal al expresar la agrupación de la suma (enlaza los resultados de $6^2$ y $8^2$ directamente a $100$), aplica correctamente la raíz cuadrada ($\sqrt{100}$) y llega a $10$ cm. \\
\textbf{Veredicto:} \textbf{Correcto} (2/2 pts).

\subsection*{3. Vectores (Fuerza resultante)}
\textbf{Procedimiento y resultado:} La estudiante es muy ordenada al desglosar el cálculo: resuelve la sumatoria en el eje $x$ ($45 - 15 = 30$) y luego en el eje $y$ ($25 + 35 = 60$). Ensambla exitosamente el vector resultante $\vec{F} = 30\hat{i} + 60\hat{j}$ N (aunque su caligrafía en los vectores unitarios finales es algo confusa, las magnitudes están correctamente situadas). \\
\textbf{Veredicto:} \textbf{Correcto} (2/2 pts).

\subsection*{4. Potencias de diez (Notación científica y prefijo)}
\textbf{Procedimiento y resultado:} Expresa la medida decimal de forma adecuada empleando notación científica estricta ($1,25 \times 10^{-5}$). Sin embargo, omite la respuesta final específica requerida en el enunciado, la cual consistía en identificar e indicar el nombre del prefijo del Sistema Internacional (micrómetros). \\
\textbf{Veredicto:} \textbf{Incompleto / Parcialmente correcto} (1/2 pts).

\subsection*{5. Sistemas de unidades (Conversión de densidad)}
\textbf{Procedimiento y resultado:} Plantea los factores de conversión en cadena de forma excelente, escribiendo las fracciones exactas para transformar la masa y el volumen. Aunque escribe un resultado intermedio extraño a la derecha de la igualdad, inmediatamente corrige/aclara el cálculo en la línea inferior escribiendo "$1,03 \times 1000 = 1030 \text{ Kg/m}^3$", lo cual es un procedimiento final y resultado totalmente correctos. \\
\textbf{Veredicto:} \textbf{Correcto} (2/2 pts).

\vspace{0.5cm}
\section*{Comentarios Generales y Sugerencias}
El examen es muy bueno. La estudiante demuestra un gran manejo procedimental, ordenando sus cálculos de modo que es fácil seguir la lógica matemática en cada respuesta. 
\begin{itemize}
    \item Se recomienda prestar atención a los detalles teóricos de las instrucciones para asegurar responder a todos los requerimientos (como ocurrió con el prefijo en la Pregunta 4).
\end{itemize}

\end{document}